\chapter{Differential graded modules and resolutions}
\label{ch:dg-modules}

Fix a commutative ring \(k\). All gradings are cohomological and may be
unbounded in both directions. For a homogeneous element \(x\) we write
\(|x|\) for its degree. The tensor differential on complexes of
\(k\)-modules is
\[
 d(x\otimes y)=dx\otimes y+(-1)^{|x|}x\otimes dy .
\]
The sign rule for this tensor product is also the sign rule for module
actions and for composition. Throughout Chapter~\ref{ch:dg-modules}
we work with right modules. See \cite[\S\S2.2, 3.1--3.2]{keller}
and \cite[\S\S3.3--3.4]{yekutieli} for the background.

\section{Actions, maps, and cones}

\begin{definition}[Differential graded algebras and categories]\label{res:G01}
A differential graded algebra \(A\) is a unital graded \(k\)-algebra
with a differential of degree \(1\) satisfying \(d^2=0\), \(d1=0\), and
\[
 d(ab)=d(a)b+(-1)^{|a|}a\,d(b).
\]
A small dg category \(\mathcal A\) has a set of objects, a complex
\(\mathcal A(a,b)\) for each pair, closed identities of degree \(0\),
and associative unital composition satisfying the same rule with
\(ab\) read as \(a\circ b\). Its category \(H^0(\mathcal A)\) has morphisms
\(H^0(\mathcal A(a,b))\). A dg functor preserves identities,
composition, and the differentials. The one-object construction
identifies dg algebras with one-object dg categories.
\end{definition}
The Leibniz rule makes composition a chain map on the tensor product.
It also shows that composing a boundary with a closed degree-zero
morphism yields a boundary. Hence \(H^0(\mathcal A)\) really is a category.

\begin{definition}[Right modules over a dg category]\label{res:G02}
A right dg \(\mathcal A\)-module \(M\) consists of complexes \(M(a)\)
with associative unital actions
\(M(b)\otimes_k\mathcal A(a,b)\to M(a)\). Write an action as \(x u\).
It satisfies
\[
        d(xu)=d(x)u+(-1)^{|x|}x\,d(u).
\]
The action axioms require \((x u)v=x(u\circ v)\) for composable arrows.
A closed degree-zero map of modules is a family of chain maps
preserving these actions. A quasi-isomorphism of modules is a map
inducing cohomology isomorphisms objectwise.
\end{definition}
The definition for a one-object \(\mathcal A\) is exactly a right dg
module over \(A\). Kernels and cokernels of the degree-zero maps are
objectwise because they are stable under the actions and differentials.

\begin{proposition}[Hom complex and homotopy]\label{res:G03}
There is a dg category of right dg modules whose morphism complex
\(\underline{\operatorname{Hom}}_{\mathcal A}(M,N)\) in degree \(r\)
consists of families of degree-\(r\) graded \(k\)-module maps
\(f_a\) satisfying \(f(xu)=f(x)u\). Its differential is
\[
        \partial f=d_N f-(-1)^r f d_M .
\]
In \(K(\mathcal A)=H^0(\mathrm{dgMod}\text{-}\mathcal A)\), two chain
maps are equal precisely when their difference is \(d_Nh+hd_M\)
for a degree-\(-1\) map \(h\) preserving the right action.
\end{proposition}
\begin{proof}
We use the sign convention with no extra sign in graded
right-linearity. On the expression for \((\partial f)(xu)\),
the two terms involving \(d(u)\) are
\((-1)^{|x|+r}f(x)d(u)\) and
\(-(-1)^{r+|x|}f(x)d(u)\). They cancel, so \(\partial f\) is
again right-linear. Substituting \(\partial f\) into the same
formula gives \(\partial^2f=0\), since \(d_N^2=d_M^2=0\) and
the middle terms have opposite signs. For homogeneous \(g\)
of degree \(s\) we have
\[
 \partial(gf)
 =d_Pgf-(-1)^{s+r}gfd_M
 =(\partial g)f+(-1)^s g(\partial f)
\]
where \(P\) is the target of \(g\). This verifies the dg category
structure. The statement about \(K(\mathcal A)\) follows by taking
the degree-zero cycles and boundaries in this very complex.
\end{proof}

\begin{proposition}[Shift, cone, and the acyclicity test]\label{res:G04}
For \(s\in\mathbb Z\), put \(M[s]^i=M^{i+s}\) with differential
\((-1)^s d_M\) and the same right action. For a closed degree-zero
map \(f:M\to N\), its cone is \(N\oplus M[1]\) with differential
\[
 d(y,x)=(d_N y+f(x),-d_M x).
\]
The map \(f\) is a quasi-isomorphism if and only if its cone is
acyclic. The cone and shift act objectwise. The homotopy category
\(K(\mathcal A)\) is triangulated with the triangles isomorphic to
these cone triangles.
\end{proposition}
\begin{proof}
After a shift the Leibniz term acquires the sign
\((-1)^{s+|x|_M}=(-1)^{|x|_{M[s]}}\). The shift thus respects the
right action. On a homogeneous map of degree \(r\), multiplication
by \((-1)^{sr}\) upon reindexing gives the dg shift functor, as
substitution into \(\partial\) verifies. The cone differential
squares to zero since \(f\) is a chain map, and the term \(f(xu)\)
equals \(f(x)u\). Thus the cone is a dg module. The short exact
sequence \(0\to N(a)\to C(f)(a)\to M(a)[1]\to 0\) identifies
its connecting homomorphism with \(H(f(a))\) by lifting a cycle
as \((0,x)\). Exactness gives the equivalence.

The cone proof of the triangulated axioms in
\cite[K01--K06]{ha-core} applies with right-linear components.
More precisely, the comparison of homotopic maps, the rotation,
the map of cones obtained from a specified homotopy of squares,
and the octahedral maps there use only identities, compositions,
sums, shifts and the cone matrix. All these formulas preserve
right-linearity by the verifications above. The component
homotopies in the cone proof are also right-linear, since they
are matrices formed from the same operations on the input
homotopies. Therefore these are triangles and the same comparison
maps inside \(K(\mathcal A)\).
\end{proof}

\section{Representables and free cells}

\begin{lemma}[Enriched Yoneda for right modules]\label{res:G05}
For \(a\in\mathcal A\) put \(h_a(b)=\mathcal A(b,a)\) with the
right action by composition. Evaluation at \(1_a\) is a natural
isomorphism of complexes
\[
     \underline{\operatorname{Hom}}_{\mathcal A}(h_a,M)
        \longrightarrow M(a).
\]
The assignment \(a\mapsto h_a\) is a fully faithful dg functor.
\end{lemma}
\begin{proof}
For \(m\in M(a)^r\), its inverse map has components
\(u\mapsto m u\) on \(\mathcal A(b,a)\). This is graded right-linear.
Any graded map out of \(h_a\) must have that formula since
\(u=1_a u\). The terms involving \(d u\) in its Hom differential
cancel by the Leibniz rule, leaving the map \(u\mapsto (d m)u\).
On representables the Yoneda map takes composition to composition
as \((m u)v=m(u\circ v)\). These checks give the asserted dg
functor and full faithfulness.
\end{proof}

A free graded cell of degree \(n\) at \(a\) is a copy of
\(h_a[-n]\) with a named element \(e\) corresponding to \(1_a\).
One can adjoin it to a dg module \(P\) by setting \(de=z\) for
a chosen cycle \(z\in P(a)^{n+1}\). On products the
differential must be
\[
                 d(eu)=z u+(-1)^n e\,d u.
\]
This formula squares to zero. Indeed \(d(z u)=(-1)^{n+1}z\,du\),
which cancels the contribution from differentiating the other
term. The underlying graded module is the direct sum with the
new cell. Its quotient by \(P\) is \(h_a[-n]\). A semi-free module
here has an exhaustive filtration
\(0=P_{-1}\subset P_0\subset P_1\subset\cdots\)
obtained by adjoining arbitrary sets of such cells, each with
differential in the preceding stage. All sums are direct sums of
graded modules, so each element involves finitely many cells.

\begin{theorem}[Unbounded semi-free resolution]\label{res:G06}
Every dg \(\mathcal A\)-module \(M\) admits a degreewise surjective
quasi-isomorphism \(P\to M\) from a semi-free module.
\end{theorem}
\begin{proof}
Put in stage zero a cell of degree \(n\) with zero differential
for every cycle in \(M(a)^n\) at every object and every degree.
Map the generator to that cycle. Suppose \(p_j:P_j\to M\) is
defined. For each \(z\in P_j(a)^n\) with \(dz=0\), and each
\(x\in M(a)^{n-1}\) satisfying \(dx=p_j(z)\), adjoin in stage \(j+1\)
a generator \(e\) of degree \(n-1\) with \(de=z\) and send it to \(x\).
The displayed cell formula and the equation on the proposed image
make this a chain map. The indexing uses only sets because
\(\mathcal A\) is small.

Any class in \(H^n(M(a))\) is represented by a cycle from stage zero.
A cycle of \(P(a)\) whose image is a boundary is already a cycle
at a finite stage because the filtrations consist of injections
and each element has finite support. At the next stage a cell
with a chosen primitive of its image kills that cycle. Hence
\(H^n(P(a))\to H^n(M(a))\) is both surjective and injective.
Finally, for an arbitrary \(x\in M(a)^n\) the cycle \(dx\) was
represented in stage zero. The new cell attached for this cycle
and \(x\) gives a degree-\(n\) preimage at stage one. This proves
degreewise surjectivity.
\end{proof}
